% =====================================================================
%  İÜC Endüstri Mühendisliği — GÜZ Dönemi Ders Çizelgeleme Problemi
%  Kısıt Programlama (CP-SAT) Matematiksel Modeli
%  Derleme: xelatex matematiksel_model.tex
% =====================================================================
\documentclass[11pt,a4paper]{article}

\usepackage{fontspec}
\setmainfont{Times New Roman}
\usepackage{polyglossia}
\setmainlanguage{turkish}

\usepackage{amsmath,amssymb,amsthm,mathtools}
\usepackage[margin=2.6cm]{geometry}
\usepackage{booktabs}
\usepackage{array}
\usepackage{enumitem}
\usepackage{xcolor}
\definecolor{iucnavy}{RGB}{18,44,84}
\definecolor{iucaccent}{RGB}{140,30,45}
\usepackage{titlesec}
\titleformat{\section}{\Large\bfseries\color{iucnavy}}{\thesection}{0.8em}{}
\titleformat{\subsection}{\large\bfseries\color{iucnavy}}{\thesubsection}{0.7em}{}
\usepackage[colorlinks=true,linkcolor=iucaccent,urlcolor=iucaccent]{hyperref}
\usepackage{microtype}

\newtheorem{proposition}{Önerme}
\newtheorem{remark}{Not}
\theoremstyle{definition}
\newtheorem{definition}{Tanım}

\newcommand{\SET}[1]{\mathcal{#1}}
\newcommand{\st}{\quad\text{öyle ki}\quad}
\DeclareMathOperator{\rank}{rank}

\begin{document}

% ------------------------------------------------------------------ başlık
\begin{center}
  {\color{iucnavy}\LARGE\bfseries Üniversite Ders Çizelgeleme Probleminin\\[2pt]
  Kısıt Programlama Modeli}\\[8pt]
  {\large Endüstri Mühendisliği Bölümü\\
  Lisans, Yüksek Lisans ve Doktora Programlarının Bütünleşik Eniyilemesi}\\[6pt]
  {\small\itshape Müfredat-tabanlı ders çizelgeleme (CB--CTT) ailesinin,
  dinamik seçmeli eşleştirme ve çok-kaynaklı oda yönetimi ile genişletilmiş bir örneği}\\[10pt]
  {\normalsize\bfseries Prof.\,Dr.\ Tarık KÜÇÜKDENİZ}\\[2pt]
  {\small İstanbul Üniversitesi--Cerrahpaşa, Endüstri Mühendisliği Bölümü}\\[6pt]
  \rule{0.6\textwidth}{0.6pt}
\end{center}

\section{Problem Tanımı}

Bir eğitim döneminde; üç program düzeyine (dönemde açık lisans yarıyılları ---
güzde 1., 3., 5., 7.; baharda 2., 4., 6., 8.--- ile yüksek lisans ve doktora) ait yaklaşık elli dersin, beş iş günü $\times$ on üç ders
saatinden oluşan haftalık ızgaraya; derslik, laboratuvar, amfi ve öğretim üyesi
kaynakları paylaşılarak, hiçbir öğrenci ve öğretim üyesi çakışması olmadan
yerleştirilmesi istenmektedir. Fakülte tarafından saatleri dışarıdan sabitlenmiş
ortak dersler (Matematik~I, Fizik~I, Genel Kimya, Programlamaya Giriş) hareket
ettirilemez engeller olarak modele girer. Problem, literatürdeki
\emph{Curriculum-Based Course Timetabling} (CB--CTT, ITC--2007 Track~3) sınıfının;
\begin{itemize}[nosep]
  \item seçmeli derslerin \emph{çözücü tarafından seçilen} eşleşmelerle aynı saate
        bindirilebilmesi (dinamik eşleştirme),
  \item ders bazında oda zorunlulukları, çift-oda rezervasyonu ve laboratuvar
        kuralları,
  \item her dönem için \emph{tam boş gün} garantisi
\end{itemize}
ile zenginleştirilmiş bir uzantısıdır. Model, Google OR-Tools \textsc{CP-SAT}
çözücüsüyle eniyilenir ve elde edilen çizelge, modelden bağımsız bir doğrulayıcı
tarafından tüm sert kısıtlara karşı yeniden denetlenir.

\section{Kümeler ve İndisler}

\begin{center}
\renewcommand{\arraystretch}{1.25}
\begin{tabular}{@{}>{$}l<{$} p{0.78\textwidth}@{}}
\toprule
\SET{D}=\{1,\dots,5\} & iş günleri (Pazartesi \dots\ Cuma) \\
\SET{P}=\{1,\dots,13\} & günlük ders saatleri (50 dakikalık periyotlar; $1\equiv$ 08.30) \\
\SET{K} & açılacak dersler kümesi \\
\SET{C} & \emph{bileşenler}: her dersin ayrık olarak çizelgelenen T (teori),
          U (uygulama), L (laboratuvar) veya birleşik blokları; $c\in\SET{C}$,
          süresi $\ell_c$ periyot \\
\SET{R} & derslikler: lisans havuzu, lisansüstü havuzu, bilgisayar laboratuvarı
          $r^{\mathrm{lab}}$ ve amfi $r^{\mathrm{amfi}}$ \\
\SET{H} & öğretim elemanları (uygulama asistanları dâhil); $h(c)\in\SET{H}$
          bileşenin yürütücüsü \\
\SET{Q} & öğrenci dönemleri (kohortlar): dönemde açık lisans yarıyılları ile YL ve DR; $q(c)\in\SET{Q}$ \\
\SET{G} & seçmeli ders grupları; $g\in\SET{G}$ için üye bileşen kümesi
          $E_g\subseteq\SET{C}$, üye sayısı $n_g=|E_g|$ ve seçim zorunluluğu
          $k_g$ (öğrenci $g$'den tam $k_g$ ders seçer) \\
\SET{F} & dışarıdan sabitlenmiş ortak oturumlar; $f=(\delta_f,\pi_f,\ell_f,r_f,[h_f])$ \\
\SET{B}\subset\SET{D}\times\SET{P} & kampüs geneli yasak dilimler
          (ör.\ Cuma namazı: $(5,6)$) \\
\bottomrule
\end{tabular}
\end{center}

\section{Parametreler}

\begin{center}
\renewcommand{\arraystretch}{1.25}
\begin{tabular}{@{}>{$}l<{$} p{0.74\textwidth}@{}}
\toprule
R_c\subseteq\SET{R} & $c$ bileşeninin \emph{uygun oda kümesi}: düzey havuzu
  (lisansüstü havuzu bilgisayar laboratuvarını da içerir; uygun lisans bilgisayar
  dersleri havuza ek olarak laboratuvarı kullanabilir); laboratuvar dersleri için
  $\{r^{\mathrm{lab}}\}$; yüksek mevcutlu dersler için $\{r^{\mathrm{amfi}}\}$;
  çift-oda dersleri derslik havuzundan seçer (laboratuvar ayrıca rezerve edilir);
  ders bazlı oda zorunlulukları (ör.\ D-731) ile kesişim \\
A_h\subseteq\SET{D}\times\SET{P} & $h$ hocasının \emph{uygun olmadığı} dilimler
  (gün/saat pencereleri) \\
{[p_{\min},\,p_{\max}]} = [2,\,10] & bölüm derslerine açık günlük pencere
  (08.30 dilimi ve 17.40 sonrası yasak) \\
\kappa = 2 & aynı dönemin aynı anda yürüyebilecek en çok seçmeli ders sayısı \\
w_{\bullet} & amaç fonksiyonu ağırlıkları (Tablo~\ref{tab:weights}) \\
\bottomrule
\end{tabular}
\end{center}

\section{Karar Değişkenleri}

\begin{align}
& d_c \in \SET{D}, \qquad
  p_c \in \{\,p_{\min},\dots,p_{\max}-\ell_c+1\,\}
  && \text{bileşenin günü ve başlangıç saati} \notag\\[2pt]
& t_c = 13\,(d_c-1) + p_c, \qquad
  I_c = [\,t_c,\; t_c+\ell_c\,)
  && \text{düzleştirilmiş zaman ve aralık} \notag\\[2pt]
& y_{c,r} \in \{0,1\}, \quad r\in R_c,
  \qquad \textstyle\sum_{r\in R_c} y_{c,r} = 1
  && \text{oda ataması} \notag\\[2pt]
& s_{uv} \in \{0,1\}, \quad \{u,v\}\subseteq E_g
  && \text{aynı-slot eşleşme göstergesi} \notag
\end{align}

\noindent Ayrıca doğrusallaştırma için türetilmiş göstergeler kullanılır:
kohort--gün kullanım değişkenleri $u_{q,\delta}$, hoca--gün kullanım değişkenleri
$v_{h,\delta}$ ve gün-içi kompaktlık büyüklükleri (ilk/son ders, meşguliyet).

\section{Sert Kısıtlar}

\subsection*{K1 --- Zaman Penceresi}
Tanım gereği $p_c\ge p_{\min}$ ve $p_c+\ell_c-1\le p_{\max}$: hiçbir bölüm dersi
08.30 diliminde başlayamaz, 17.35'ten sonra süremez ve gün sınırını aşamaz.
Kampüs yasak dilimleri ile hoca uygunlukları başlangıç bölgesinden budanır:
\begin{equation}
\forall c,\;\forall(\delta,\pi)\in A_{h(c)}\cup\SET{B}:\qquad
\neg\Bigl[\,d_c=\delta \;\wedge\; p_c \le \pi \le p_c+\ell_c-1\,\Bigr].
\end{equation}

\subsection*{K2 --- Oda Tek-Doluluk}
Her oda her an en çok bir oturum barındırır; sabit ortak oturumlar ve dış
rezervasyonlar da odalarını tüketir:
\begin{equation}
\forall r\in\SET{R}:\quad
\textsc{NoOverlap}\Bigl(\{\,I_c : y_{c,r}=1\,\}\;\cup\;\{\,I_f : r_f=r\,\}\Bigr).
\end{equation}
\emph{Çift-oda dersleri} (ilk haftalar sınıf, sonrasında laboratuvar) hem seçilen
dersliği hem laboratuvarı aynı anda rezerve eder: $c\in\text{Dual}$ için $I_c$,
$r^{\mathrm{lab}}$'ın NoOverlap kümesine koşulsuz eklenir.

\subsection*{K3 --- Öğretim Üyesi Tek-Doluluk ve Uygunluk}
\begin{equation}
\forall h\in\SET{H}:\quad
\textsc{NoOverlap}\Bigl(\{\,I_c : h(c)=h\,\}\;\cup\;\{\,I_f : h_f=h\,\}\Bigr),
\end{equation}
sabitlenmiş ortak dersi olan hocanın (ör.\ programlama laboratuvarı) o saatleri
de hesaba katılır. Uygunluk pencereleri K1'de işlenmiştir.

\subsection*{K4 --- Öğrenci Çakışmasızlığı (dönem içi)}
$Z_q$, $q$ kohortunun zorunlu bileşenleri; $S_q$ seçmelileri olsun.
\begin{align}
&\textsc{NoOverlap}(Z_q) && \text{zorunlular ikili ayrık,}\\
&\forall e\in S_q:\;\; \textsc{NoOverlap}\bigl(Z_q\cup\{e\}\bigr)
  && \text{her seçmeli, zorunlu iskeletle ayrık,}\\
&\forall g\ne g' ,\; u\in E_g,\, v\in E_{g'}:\;\; I_u\cap I_v=\emptyset
  && \text{farklı gruplardan dersler ayrık}
\end{align}
(son kural, öğrencinin her iki gruptan da ders aldığı 7.~yarıyıl F--G çiftini
kapsar). Ayrıca sabit ortak oturumun bağlı olduğu kohortun \emph{tüm} bileşenleri
o oturumla zaman-ayrıktır: $q(c)=q_f \Rightarrow I_c\cap I_f=\emptyset$.%
\footnote{K4 lisans kohortları için kurulur; mevcut lisansüstü kohortlarda
çizelgelenen zorunlu bileşen ve seçmeli grubu bulunmadığından bu kısıtlar boş
sağlanır; lisansüstü seçmelileri yalnızca K6 ($\kappa=2$) ile sınırlıdır.}

\subsection*{K5 --- Dinamik Aynı-Slot Eşleştirme}
Aynı grubun herhangi iki dersi \emph{ya birebir aynı slotta ya tamamen ayrık}
olabilir; kısmi örtüşme yasaktır. Eşleştirmeyi konfigürasyon değil,
\emph{çözücü} seçer:
\begin{align}
s_{uv}=1 &\;\Longrightarrow\; d_u=d_v \,\wedge\, p_u=p_v, \\
s_{uv}=0 &\;\Longrightarrow\; t_u+\ell_u\le t_v \;\vee\; t_v+\ell_v\le t_u,
\qquad \{u,v\}\subseteq E_g,\; k_g\ge 2,\\
&\sum_{\{u,v\}\subseteq E_g} s_{uv} \;\le\; \max\bigl(0,\; n_g-k_g\bigr)
\qquad (k_g \ge 2). \label{eq:cap}
\end{align}
Seçim-1 gruplarında ($k_g=1$) üyeler serbestçe örtüşebilir; K6 kapasitesi yine
geçerlidir.

\begin{proposition}[Seçim Garantisi]\label{prop:selection}
Eş uzunluklu üyeli bir $g$ grubunda \textup{(\ref{eq:cap})} ve K6 kapasitesi
$\kappa=2$ altında, üyelerin başlangıç zamanlarına göre denklik sınıfları en az
$k_g$ farklı slot oluşturur ve her sınıftan en çok bir ders seçen \emph{her}
öğrenci, $k_g$ dersini ikili çakışmasız alır.
\end{proposition}
\begin{proof}[İspat taslağı]
Kısmi örtüşme yasak olduğundan ``aynı slot'' ilişkisi, eş uzunluk altında bir
denklik bağıntısıdır; $\kappa=2$ her sınıfın en çok iki üye içermesini sağlar.
Böylece sınıf sayısı $n_g-\sum s_{uv}\ge n_g-(n_g-k_g)=k_g$'dir. Farklı
sınıflardan dersler $s_{uv}=0$ gereği ayrık aralıklarda yattığından, sınıf
başına en çok bir ders içeren her $k_g$-altküme ikili ayrıktır. \qedhere
\end{proof}

\subsection*{K6 --- Eşzamanlı Seçmeli Kapasitesi}
Aynı kohortun seçmelileri için birim talepli kümülatif kısıt:
\begin{equation}
\forall q:\quad \textsc{Cumulative}\bigl(\{I_e : e\in S_q\},\;
\mathbf{1},\; \kappa\bigr),\qquad \kappa=2 .
\end{equation}

\subsection*{K7 --- Yapısal Ders Kuralları}
\begin{itemize}[nosep]
\item \emph{Ardışık bloklarda tek oda}: Bir dersin laboratuvar dışı iki bileşeni
  zamanda \emph{ardışık} ise (biri bittiği an diğeri başlıyorsa) aynı dersliği
  kullanır; kesintisiz bir ders akışı oda değiştirmez. Ardışıklık reifiye bir
  Boole ile algılanır ve oda eşitliğini koşullu tetikler:
  $a_{uv}\Rightarrow y_{u,r}=y_{v,r}$. Ardışık olmayan bloklar (farklı gün veya
  araya boşluk giren sabah/öğleden sonra blokları) farklı derslik kullanabilir.
\item \emph{Tek blok}: İşaretli dersler tüm saatleriyle tek kesintisiz bileşen
  olarak üretilir (ör.\ programlama dersi $T{+}L=4$ saat).
\item \emph{Aynı gün}: Belirlenen ders çiftlerinin bileşenleri ortak gün alır:
  $d_u=d_v$.
\item \emph{Farklı gün}: Bölünmüş derslerin parçaları farklı günlere düşer:
  $d_u\ne d_v$ (ör.\ tesis planlama $2{+}2$).
\item \emph{Gün kısıtları}: Belirli dersler yapılandırmayla belirli günlere
  sınırlanır ($d_c\in\SET{D}_c$; ör.\ üç lisans dersi $\to$ Salı); istenirse bir
  kohortun tamamı da bir gün kümesine sınırlanabilir.
\item \emph{Sabitlemeler}: Ders bazında gün/saat(/oda) pinleri $d_c,p_c,(y_{c,r})$
  değerlerini eşitlikle bağlar; eşleşmeyen pin girdileri model kurulurken hata
  üretir.
\end{itemize}

\subsection*{K8 --- Tam Boş Gün Garantisi}
$u_{q,\delta}\in\{0,1\}$; $q$ kohortunun (zorunlu \emph{ve tüm} seçmeli
bileşenleri ile o kohorta bağlı sabit ortak oturumları dâhil) $\delta$ gününü
kullanıp kullanmadığını reifikasyonla izler; sabit oturum günlerinde
$u_{q,\delta}=1$'e sabitlenir. Garanti:
\begin{equation}
\forall q\in\SET{Q}:\qquad \sum_{\delta\in\SET{D}} u_{q,\delta} \;\le\; 4 .
\end{equation}
Böylece her dönemin öğrencisi, seçmeli tercihinden bağımsız olarak haftada en az
bir tam gün derssizdir.

\section{Amaç Fonksiyonu}

Sözlükbilimsel öncelikleri yaklaşıklayan ağırlıklı tek ölçüt en küçüklenir:
\begin{equation}
\min\;
\underbrace{w_1\!\!\sum_{q,\delta} u_{q,\delta}}_{\text{boş gün}}
+\underbrace{w_2\!\!\sum_{h,\delta} v_{h,\delta}}_{\text{hoca-gün}}
+\underbrace{w_3\!\!\sum_{q,\;\delta\notin \Pi_q} u_{q,\delta}}_{\text{gün tercihi}}
+\underbrace{w_4\!\!\sum_{c} m_c}_{\text{Pzt.\ erken}}
+\underbrace{w_5\!\!\sum_{q,\delta} \gamma_{q,\delta}}_{\text{boşluk}}
+\underbrace{w_6\!\!\sum_{g}\sum_{uv} s_{uv}}_{\text{eşleşme}}
+\underbrace{w_7\!\!\sum_{c,r} \rho_r\, y_{c,r}}_{\text{oda tercihi}}
\end{equation}
Burada $v_{h,\delta}$ hoca--gün kullanımı (sabit oturum günleri dâhil);
$\Pi_q$ kohortun tercih edilen gün kümesi (ör.\ L5 için Pzt--Çar);
$m_c$, Pazartesi 09.25 başlangıcının göstergesi;
$\gamma_{q,\delta}=\max\bigl(0,\ \text{süre}(q,\delta)-\text{meşguliyet}(q,\delta)\bigr)$
temsilî öğrenci yüküyle hesaplanan gün-içi boşluk (negatif kırpma, meşru
örtüşmelerin amaç mekanizmasınca yasaklanmamasını sağlar);
$\rho_r$ düzeye özgü oda öncelik sırası (lisans: D-731 önce; lisansüstü:
Lisansüstü Sınıfı $\succ$ D-724/D-733 $\succ$ TUSAŞ).%
\footnote{Yapılandırmada ayrıca $w_{\mathrm{P1}}=25$ yedek ağırlığı tanımlıdır;
K1'in sert 08.30 yasağı altında ilgili terim özdeş olarak sıfırdır ve yalnızca
yasak kaldırılırsa etkinleşir; bu nedenle amaç ifadesinde gösterilmemiştir.}

\begin{table}[h]
\centering
\caption{Ağırlık hiyerarşisi: üst katmanlar alt katmanlara pratikte baskındır
(yaklaşık sözlükbilimsel sıralama).}\label{tab:weights}
\renewcommand{\arraystretch}{1.15}
\begin{tabular}{@{}llr@{}}
\toprule
Katman & Ölçüt & Ağırlık \\
\midrule
Birincil & öğrenci boş günü ($w_1$) & 1000 \\
Birincil & hoca gün yoğunlaştırma ($w_2$) & 500 \\
İkincil & kohort gün tercihi ($w_3$) & 300 \\
İkincil & Pazartesi 10.20 başlangıcı ($w_4$) & 150 \\
Üçüncül & gün-içi öğrenci boşluğu ($w_5$) & 10 \\
Kırılım & eşleşme cezası ($w_6$) & 5 \\
Kırılım & oda önceliği ($w_7$) & 2 \\
\bottomrule
\end{tabular}
\end{table}

\begin{remark}
Eşleşme cezası $w_6$ bilinçli olarak \emph{küçüktür}: seçmelileri üst üste
bindirmek başlı başına bir hedef değildir; yalnızca daha üst katmanlarda gerçek
bir kazanım (boş gün, gün tercihi, fizibilite) sağlıyorsa gerçekleşir.
Kararsızlık durumunda dersler ayrık kalır ve öğrenci seçim özgürlüğü en geniş
hâlinde korunur.
\end{remark}

\section{Çözüm Yöntemi ve Doğrulama}

Model, aralık değişkenleri ile $\textsc{NoOverlap}$/$\textsc{Cumulative}$ küresel
kısıtları üzerine kurulmuş bir \textsc{CP-SAT} örneğidir: $|\SET{C}|=48$ bileşen,
$5\times 9$ etkin zaman ızgarası ve 11 oda üzerinde \emph{birincil} karar
değişkenleri 96 gün/başlangıç tamsayısı ($d_c,p_c$), 150 oda-atama Boole'u
($y_{c,r}$) ve 36 eşleşme Boole'undan ($s_{uv}$) oluşur. Reifikasyon ve amaç
doğrusallaştırması ile birlikte tam model \textbf{2\,116 karar değişkeni}
(1\,129 Boole, 987 tamsayı) ve \textbf{4\,815 kısıt} içerir; yalnız sert
çekirdek 1\,918 değişken / 4\,539 kısıttır. Deterministik aramayla
dakikalar içinde \emph{kanıtlı en iyi} (OPTIMAL) çözüm elde edilmektedir. Çözüm sonrası,
modelden bağımsız bir denetleyici; oda/hoca çakışmaları, pencere ihlalleri,
grup--slot yapısı (Önerme~\ref{prop:selection} koşulları dâhil), sabit oturum
çakışmaları ve boş gün garantisini ham çizelge üzerinden yeniden hesaplayarak
sıfır ihlal raporlar. Tüm veri girdileri (uygunluklar, gruplar, oda kuralları,
ağırlıklar) koda dokunmadan yapılandırma dosyalarından yönetilir.

\vspace{6pt}
\begin{center}
\rule{0.6\textwidth}{0.6pt}\\[4pt]
{\small\itshape Bu belge, uygulamadaki modelin birebir matematiksel karşılığıdır;
her kısıt ailesi kaynak kodda ayrı birim testleriyle korunmaktadır.}
\end{center}

\end{document}
